\subsection{Integrating the crossed rational subgroup first}
\label{sec:comp-crossed-first}
Before integrating the crossed box in the rationalized angle, separate
the scalar terms whose coefficients and function arguments are rational
in the original invariant $y$.
Transforming those terms to the angle variable would introduce an auxiliary root $\rho_u=\sqrt{(x-1)/(x+3)}$, whose
apparent threshold at $x=1$ is absent from the original scalar
integral. This root is introduced by applying the spacelike
variable $v$ to a subgroup that is already rational in the
original invariant $y$.
This subgroup can be integrated before that change of variable.
Here $x=s-1>0$ and $x/(1+x)\leq y\leq x$, in the mass units
used for the box calculation.

For this subgroup write $A=\overline A_1$, $B_s=\overline B_{02}$,
and $C_s=C_{012}$, evaluated at the crossed invariant $S=1-y$.
The subscript $s$ in these local scalar names does not label the
photon ordering. Include also the explicit constant in the crossed-box
trace. The scalar functions are
\begin{equation}
 A=1,\qquad B_s=2+\frac{y\ln y}{1-y},\qquad
 C_s=\frac{\operatorname{Li}_2(1-y)-\pi^2/6}{y}.
\end{equation}
Define the two GPLs
\begin{equation}
 F(y)=G(1,0;y)=\frac{\pi^2}{6}-\operatorname{Li}_2(1-y),
 \qquad W(y)=G(0,1,0;y).
\end{equation}
Their differential equations are
$F'(y)=\ln y/(y-1)$ and $W'(y)=F(y)/y$.
In particular, $F$ is real and analytic for every $y>0$.
The apparent pole of its derivative at $y=1$ is removable.
These statements fix the continuation of the GPL combination
without assigning separate imaginary parts at an artificial pole.

The selected crossed-box integrand is
\begin{equation}
 I_{\mathrm{rat}}(x,y)=R(x,y)+Q(x,y)\ln y-D(x,y)F(y),
\end{equation}
where exact partial fractions of the trace coefficients give
\begin{align}
 R={}&\frac{16(x-2)}{y^2}+\frac{28x+12}{x}
       +\frac4{x(y-1)}-\frac{8(2x^2+x-4)}{xy},\\
 Q={}&\frac{8(x-2)}y+\frac{4(x-6)}{x(y-1)}
       -\frac{12x+4}{x}-\frac4{x(y-1)^2},\\
 D={}&-\frac{8y}x+\frac{8x-8}x-\frac{4(x^2-8)}{xy}.
\end{align}
The only apparent double pole is integrated with its accompanying
simple pole. The exact identity
\begin{equation}
 \frac1{y-1}-\frac{\ln y}{(y-1)^2}
 =\frac{\dd}{\dd y}\frac{\ln y}{y-1}+\frac1y
\end{equation}
shows explicitly why no isolated $\ln(1-y)$ is needed.
Integration by parts of the two polynomial coefficients of $F$
then gives the finite primitive
\begin{align}
 P(x,y)={}&-\frac{16(x-2)}y+\frac{(32x+28)y+y^2}{x}\\
 &+\frac{-16x^2-8x+36-(4x+16)y-2y^2}{x}\ln y
       +4(x-2)\ln^2y\\
 &+\frac4x\frac{\ln y}{y-1}
       +\frac4x[y^2+(2-2x)y+3x-9]F(y)\\
 &+\frac{4(x^2-8)}xW(y).
\end{align}
Direct differentiation, using the two displayed GPL differential
equations, gives $\partial_yP=I_{\mathrm{rat}}$ as a rational
identity in the independent symbols $\ln y,F,W$.
The contribution to the unnormalized crossed angular integral is
therefore exactly
\begin{equation}
 J_{u,\mathrm{rat}}(x)=P(x,x)-P\left(x,\frac{x}{1+x}\right).
\end{equation}
Since $\dd c=2s\,\dd y/x^2$ and the interference contributes
$(4s)^{-1}\int\operatorname{Re}\mathcal T\,\dd c$, this subgroup
enters the cross-section coefficient as $J_{u,\mathrm{rat}}/(2x^2)$.
At $x=1$ the upper endpoint quotient is assigned
its analytic value $\ln y/(y-1)=1$, so the same expression covers
the apparent root collision there.

The lower endpoint introduces no new alphabet. Write
$f(x)=x/(1+x)$. The first identity below follows directly from
the logarithm of this quotient. For the other two, differentiation
and the boundary values $F(0)=W(0)=0$ prove
\begin{align}
 \ln f(x)&=G(0;x)-G(-1;x),\\
 F(f(x))&=-G(-1,0;x)+G(-1,-1;x),\\
 W(f(x))&=-G(0,-1,0;x)+G(0,-1,-1;x)\\
        &\quad+G(-1,-1,0;x)-G(-1,-1,-1;x).
\end{align}
Thus this entire subgroup has constant letters $\{-1,0,1\}$.
The remaining crossed masters are
$B_a=\overline B_{01}$, $B_t=\overline B_{13}$,
$C_t=C_{013}$, $C_m=C_{123}$, and $D_b=D_{0123}$.
Their rational trace coefficients have denominators built only
from $x$, $y$, and $x-y+4$. In particular none has a pole at
$y=1$. After the spacelike substitution their words involve
$0,\pm1,\pm\sqrt{x/(x+4)}$ and do not involve $\rho_u$.
Consequently $\rho_u$ disappears within the crossed box itself,
before adding soft or real-radiation contributions. This removes an auxiliary root within the virtual calculation.

There is also no remaining positive interior letter in the angular
remainder. If $V=x/(x+2)$ and $z=\sqrt{x/(x+4)}$, then
\begin{equation}
 z^2-V^2=\frac{4x}{(x+4)(x+2)^2}>0.
\end{equation}
Thus $0<V<z<1$. The remainder can be integrated along the real
$v$ interval without detouring around an artificial crossed root.

To fix the lower boundary in the subsequent function-basis
conversion, the series at positive $x\to0$ follows from
\begin{align}
 F(y)&=\sum_{n=1}^{\infty}y^n
             \left(\frac1{n^2}-\frac{\ln y}{n}\right),\\
 W(y)&=\sum_{n=1}^{\infty}y^n
             \left(\frac2{n^3}-\frac{\ln y}{n^2}\right).
\end{align}
They give
\begin{align}
 J_{u,\mathrm{rat}}(x)
 ={}&x(32\ln x-32)+x^2\left(\frac{92}{3}-24\ln x\right)\\
 &+x^3\left(\frac{88}{3}\ln x-\frac{236}{9}\right)
             +O(x^4\ln x).
\end{align}
In particular this unnormalized subgroup vanishes at the lower
boundary. The series fixes the boundary constant from the derived
integrand itself, without numerical fitting or a literature value.
